\section{Introduction}

Falls can resemble sitting, kneeling, or lying down in a single frame. Their distinction depends on the descent and the posture that follows. Camera-based monitoring captures this motion without a wearable device, but furniture, lighting changes, and tracking failures can interrupt observations. Skeleton sequences describe body movement compactly while reducing the visual detail passed to the classifier. Their usefulness depends on temporal modeling and the reliability of estimated joints across successive video frames.

The intended setting is an indoor assisted-living space, where caregivers review suspected falls and acknowledge events through a monitoring interface. The design objectives are to process video near the camera, continue detection during network outages, and preserve event times when records reach the cloud after reconnection.

Recent work combines temporal operators for skeleton-based fall detection. Yu et al.~\cite{yu2025tcnte} proposed TCNTE, which couples dilated temporal convolutions with a Transformer encoder and uses weighted focal loss to address class imbalance. Their implementation integrates pose estimation and tracking on a Jetson edge device. Kim et al.~\cite{kim2025yosap} developed YOSAP-LSTM for construction scenes with occlusion. It combines person detection and tracking with AlphaPose and a one-dimensional CNN-LSTM classifier, with evaluation on an embedded platform. Both studies connect motion classification with the cost and observation errors of pose extraction.

Zhang et al.~\cite{zhang2025fallmamba} introduced Fall-Mamba, which uses cross-attention to fuse video and audio features, together with multihead temporal attention and frame masking. Its masked sequence modeling motivates studying temporal interruptions in pose sequences. Chun et al.~\cite{chun2025stonegat} proposed StoneGAT, a graph attention model for individual poses that uses bone-related edge attributes, including confidence information. Its PointOut strategy drops node features during training to study resilience to missing joints. Frame masking removes observations at selected times, whereas node removal reduces the body geometry available within an observation. These distinct losses motivate separate evaluation of temporal and spatial disruption.

We propose MaskedBiMamba to combine bidirectional state-space modeling with explicit pose-quality information. Forward and backward Mamba branches~\cite{gu2023mamba} relate each posture to surrounding motion within a buffered window; temporal attention connects observations, and a learned quality gate aggregates frame features. Training-time masking exposes the classifier to interrupted sequences. This design is evaluated against convolutional, attention-based, and unidirectional state-space alternatives. Component ablations examine attention, bidirectionality, masking, and quality pooling, while controlled perturbations separate frame loss, joint loss, and confidence noise. Window, video, and event measures describe classification and the false alarms produced by temporal aggregation. The system contribution embeds this classifier in a pose-based edge-cloud prototype, pairing on-device pose extraction with Hailo acceleration and an authenticated monitoring interface. Together, the architecture, experimental analysis, and operational alert workflow examine how a compact fall detector behaves under incomplete observations and changes in recording conditions.
